\answer{ex:09:1}{(a)~$P_\infty=0.2$, 기억 $20\,\text{s}$. (b)~$P_\infty=1$, 기억 $1\,\text{s}$. (c)~기억은 잡음 진폭에 비례한다. 여섯 배.}
\answer{ex:09:2}{$P(t)=P_\infty\coth(t\sqrt{q/R})$. (a)~$\frac12\sqrt{R/q}=5\,\text{s}$, 기억의 절반이다. (b)~$t=\frac12\ln21\,\sqrt{R/q}\approx15\,\text{s}$. 높아진 \nt{ACET}가 이만큼 유지되어야 한다.}
\answer{ex:09:3}{분산 $qT/2+R/(2T)$; $T^*=\sqrt{R/q}$, 분산 $\sqrt{qR}=P_\infty$로 필터~\eqref{eq:kalman}와 같다. $(\kappa+1/\kappa)/2$배 증가: $1.25$와 $5.05$.}
\answer{ex:09:4}{$a>a_w=1-\theta/(mw)$: $w=0.041$에서 $0.15$, $w=0.045$에서 $0.22$.}
\answer{ex:09:5}{오류는 $M\approx40$--$60$에서 나타나며, $M=80$에서 끼어드는 뉴런 $1.9$개, $M=100$에서 패턴 비율 $0.86$. 다른 패턴에서 오는 간섭의 분산은 $\approx(M-1)p(1-p)/N$이므로 $M_{\max}\approx80$, $M_{\max}\propto N/p$.}
\answer{ex:09:6}{예를 들어 $\eta(a)=\eta\min\bigl(1,\max(0,(a-0.3)/0.6)\bigr)$. $\eta a$에서는 남의 뉴런 세 개가 모두 패턴에 들어갔고, $\eta(a)$에서는 하나도 들어가지 않았다. $a\ge0.9$에서의 쓰기는 같다(비율 $1.00$).}
\answer{ex:09:7}{(a)~재생 $40$번; 연습문제~\ref{ex:09:6}의 $\eta(a)$로는 영원히 불가능. (b)~재생 $200$번, $10$밤.}
